% ============================================================
%  Silver Price Forecasting — VMD–Entropy–LASSO Framework
%  Overleaf-ready LaTeX
%  Upload this .tex file plus all figures to one Overleaf project.
%  Required figure files (place in same folder as .tex):
%    fig1_pub_trend.png  fig2_lstm_architecture.png
%    fig3_research_framework.png  fig4_silver_price_split.png
%    fig5_wordcloud_silver.png    fig6_sentiment_scores.png
%    fig7_imf_decomposition.png   fig8_approximate_entropy.png
%    fig9_single_model_forecasts.png  fig10_error_barplots.png
%    fig11_interval_forecasts.png fig12_trading_strategy_illustration.png
%    fig13_trading_evaluation.png
% ============================================================

\documentclass[12pt,a4paper]{article}

% ── Packages ────────────────────────────────────────────────
\usepackage[top=2.8cm,bottom=2.5cm,left=2.8cm,right=2.8cm]{geometry}
\usepackage[T1]{fontenc}
\usepackage[utf8]{inputenc}
\usepackage[activate=false]{microtype}

\usepackage{amsmath,amssymb,mathtools,bm}
\usepackage{booktabs,array,multirow,tabularx,longtable}
\usepackage{xcolor,colortbl}
\usepackage{graphicx,caption}
\usepackage{fancyhdr,titlesec,parskip,setspace,enumitem}
\usepackage{tcolorbox}
\usepackage[colorlinks=true,linkcolor=ElsBlue,citecolor=ElsBlue,
            urlcolor=ElsBlue]{hyperref}
\usepackage{ragged2e}

% ── Colours ─────────────────────────────────────────────────
\definecolor{ElsBlue}{RGB}{0,84,166}
\definecolor{HdrBg}{RGB}{31,73,125}
\definecolor{RowA}{RGB}{235,242,250}
\definecolor{PropBg}{RGB}{209,231,209}
\definecolor{InsightBg}{RGB}{255,248,220}
\definecolor{InsightBd}{RGB}{190,140,0}
\definecolor{GrayHd}{RGB}{110,110,110}

% ── Section styles ──────────────────────────────────────────
\titleformat{\section}
  {\normalfont\large\bfseries\color{ElsBlue}}{\thesection.}{0.5em}{}
\titleformat{\subsection}
  {\normalfont\normalsize\bfseries\color{ElsBlue}}{\thesubsection}{0.5em}{}
\titleformat{\subsubsection}
  {\normalfont\normalsize\bfseries}{\thesubsubsection}{0.5em}{}

% ── Header / footer ─────────────────────────────────────────
\pagestyle{fancy}\fancyhf{}
\fancyhead[L]{\footnotesize\color{GrayHd}\textit{Silver Forecasting -- VMD--Entropy--LASSO}}
\fancyhead[R]{\footnotesize\color{GrayHd}\textit{MCX India}}
\fancyfoot[L]{\footnotesize\color{GrayHd}Preprint $\cdot$ 2025}
\fancyfoot[R]{\footnotesize\color{GrayHd}Page~\thepage}
\renewcommand{\headrulewidth}{0.4pt}
\renewcommand{\footrulewidth}{0.3pt}
\setlength{\headheight}{15pt}

% ── Table column types ───────────────────────────────────────
\newcolumntype{L}[1]{>{\RaggedRight\arraybackslash}p{#1}}
\newcolumntype{C}[1]{>{\centering\arraybackslash}p{#1}}
\newcolumntype{R}[1]{>{\RaggedLeft\arraybackslash}p{#1}}

% ── Insight box ──────────────────────────────────────────────
\tcbuselibrary{skins,breakable}
\newtcolorbox{insightbox}[1][]{
  colback=InsightBg,colframe=InsightBd,arc=3pt,boxrule=0.8pt,
  left=8pt,right=8pt,top=6pt,bottom=6pt,
  fonttitle=\bfseries\small\color{InsightBd},#1}

% ── Shorthand commands ───────────────────────────────────────
\newcommand{\cmark}{\ensuremath{\checkmark}}
\newcommand{\hdrrow}[1]{\textcolor{white}{\textbf{#1}}}

% Helper: bold + green highlight for Proposed row
\newcommand{\proprow}{\rowcolor{PropBg}}

% ── Inline figure/table helpers (no floating) ────────────────
% Tables and figures appear exactly where called.
% Use \vspace{4pt} before caption, \vspace{8pt} after.

% ============================================================
\begin{document}
\setlength{\parindent}{0pt}\setlength{\parskip}{6pt}

% ── TITLE PAGE ───────────────────────────────────────────────
\thispagestyle{empty}
\begin{center}
{\LARGE\bfseries\color{ElsBlue}
  Silver Price Forecasting and Algorithmic Trading:\\[4pt]
  A Multimodal VMD--Entropy--LASSO Framework\\[4pt]
  for the Indian Commodity Market}\\[18pt]
\hrule height 0.8pt\vspace{6pt}
{\small MCX Silver Futures (INR/kg) $\cdot$ Weekly 2000--2026
 $\cdot$ $N=951$ Observations $\cdot$ 80/20 Train--Test\\[3pt]
 Adapted from Liu et al.\ (2025) $\cdot$ Applied to the Indian Commodity Market}
\vspace{6pt}\hrule height 0.8pt
\end{center}
\vspace{8pt}

% ── HIGHLIGHTS ───────────────────────────────────────────────
\noindent{\large\bfseries\color{ElsBlue}Highlights}
\vspace{3pt}\hrule height 0.4pt\vspace{5pt}
\begin{itemize}[leftmargin=1.5em,itemsep=2pt]
  \item First application of the VMD--ApEn--LASSO framework to MCX Silver (INR/kg),
        incorporating India-specific features: USD/INR, India VIX, MCX Gold, India EPU,
        and a 12-keyword Google Trends composite index.
  \item Adaptive complexity-aware model assignment: ARIMA for 2 low-complexity IMFs
        (86.1\% of total variance) and LSTM for 7 high-complexity IMFs, guided by
        Approximate Entropy thresholding.
  \item Proposed model achieves RMSE of 63,709~INR/kg (best among decomposition
        models), outperforming VMD-ARIMA (64,510) and CEEMDAN-ARIMA (64,372);
        DM test confirms superiority at $p<0.01$ ($DM=5.89^{***}$) over VMD-ARIMA.
  \item Interval-constrained trading strategy yields 187.2\% cumulative return
        (Sharpe 1.759, max~DD 27.7\%) under Scheme~1; Scheme~1$'$ reduces
        drawdown to 10.9\% with 30 filtered trades (Sharpe 0.701).
  \item India EPU selected in 8/9 IMFs via LASSO, confirming domestic policy
        uncertainty as the dominant cross-frequency driver of MCX silver prices.
\end{itemize}

\vspace{6pt}\noindent\rule{\linewidth}{0.4pt}

% ── ABSTRACT ─────────────────────────────────────────────────
\vspace{4pt}\noindent{\large\bfseries\color{ElsBlue}Abstract}\vspace{3pt}

\noindent
Silver occupies a unique dual role as a safe-haven precious metal and a critical
industrial commodity in the Indian market. This paper adapts the multimodal
decomposition-ensemble framework of Liu et al.\ (2025) to MCX Silver futures (INR/kg)
over 951 weekly observations from 2000 to 2026. The framework sequentially applies
Variational Mode Decomposition (VMD) to extract nine Intrinsic Mode Functions (IMFs),
Approximate Entropy (ApEn) to classify each IMF by signal complexity, LASSO regression
for adaptive per-IMF feature selection, and a hybrid ARIMA/LSTM ensemble for point
forecasting. An interval-constrained trading strategy integrates point predictions with
$\pm3\%$ interval bounds to suppress trades during high-uncertainty weeks. On the
108-week test set (March 2024 -- March 2026), the Proposed model achieves RMSE of
63,709~INR/kg, the lowest among all decomposition-based models. The Diebold--Mariano
test confirms $p<0.01$ superiority over VMD-ARIMA ($DM=5.89^{***}$). In trading
simulation, Scheme~1$'$ achieves 15.2\% cumulative return with maximum drawdown of
10.9\% and annualised Sharpe of 0.701 across 30 filtered trades, demonstrating
principled capital preservation during the 2024--2026 structural regime change.

\vspace{5pt}\noindent\rule{\linewidth}{0.4pt}

% ── KEYWORDS ─────────────────────────────────────────────────
\vspace{4pt}\noindent{\large\bfseries\color{ElsBlue}Keywords}\vspace{3pt}

\noindent
Silver price forecasting; Variational Mode Decomposition (VMD); Approximate Entropy;
LASSO feature selection; LSTM; ARIMA; Indian commodity market; MCX; Google Trends;
India EPU; Interval-constrained trading strategy; Diebold--Mariano test

\newpage

% ============================================================
\section{Introduction}
% ============================================================

Silver is among the most actively traded commodities on the Multi Commodity Exchange
of India (MCX), serving simultaneously as a safe-haven asset, an industrial input for
electronics and photovoltaics, and a cultural store of value tied to festival and
wedding demand. Its price (INR/kg) is shaped by global forces---international gold
prices, Brent crude, US dollar strength---and distinctly domestic factors: USD/INR
exchange rate, Nifty~50 equity sentiment, India VIX, import policy uncertainty (India
EPU Index), and retail investor attention (Google Trends). This multi-factor,
nonlinear, non-stationary structure makes silver price forecasting a challenging open
problem with direct implications for hedgers, portfolio managers, and commodity trading
firms.

Decomposition-ensemble methods have become the leading paradigm for commodity price
forecasting. Variational Mode Decomposition (VMD) outperforms EMD-based variants by
reducing mode mixing through variational optimisation. Liu et al.\ (2025) advanced this
paradigm for Chinese carbon markets by combining VMD with Approximate Entropy complexity
classification, LASSO feature selection, and adaptive ARIMA/LSTM assignment, extended
to an interval-constrained trading strategy. Figure~1 shows the growing academic
interest in this domain.

\vspace{4pt}
\noindent\includegraphics[width=\linewidth]{fig1_pub_trend.png}
\captionof{figure}{The time trend of publication numbers in the domain of silver price
forecasting. Both WOS and Scopus databases show a marked acceleration from 2020,
motivating the development of more sophisticated forecasting methodologies.}
\label{fig:pubtrend}
\vspace{8pt}

This study is the first to apply the VMD--ApEn--LASSO pipeline to the Indian silver
market, which presents unique challenges: a $42\times$ slope acceleration in the
2024--2026 test period, India-specific macroeconomic drivers absent from prior
commodity forecasting, and the dual precious-metal/industrial demand structure of
silver.

\medskip\noindent\textbf{Contributions.}
\begin{enumerate}[leftmargin=*,itemsep=2pt]
  \item First application of VMD--ApEn--LASSO to MCX Silver (INR/kg) with a fully
        India-specific feature set: USD/INR, India VIX, MCX Gold, India EPU, and a
        12-keyword Google Trends composite.
  \item Empirical proof that ApEn-guided model assignment is essential: VMD-LSTM is
        significantly \emph{worse} than VMD-ARIMA ($DM=-14.11^{***}$), confirming
        that LSTM must not be applied to smooth low-entropy IMFs.
  \item Frequency-dependent feature revelation: gold co-movement and USD/INR drive
        long-run IMFs; India VIX, Google Trends, and EPU drive short-term oscillations.
        LASSO selects EPU in 8/9 IMFs.
  \item An interval-constrained trading strategy that reduces drawdown from 27.7\% to
        10.9\% ($-60.7\%$) while cutting trade count by 70\%.
\end{enumerate}

The paper is organised as follows: Section~2 presents data and methodology;
Section~3 reports empirical results; Section~4 concludes.

% ============================================================
\section{Data and Methodology}
% ============================================================

\subsection{Literature Review}

Table~1 summarises the most related study on which this work is based.

\vspace{4pt}
\noindent\textbf{Table 1.}~The source study adapted in this work.
\vspace{3pt}

\renewcommand{\arraystretch}{1.18}

\noindent\small
\begin{tabular}{L{2.55cm} C{0.6cm} L{2.7cm} L{1.9cm} L{2.4cm} L{2.0cm}}
\rowcolor{HdrBg}\hdrrow{Author(s)} & \hdrrow{Year} & \hdrrow{Journal} & \hdrrow{Asset} & \hdrrow{Method} & \hdrrow{Metrics} \\
\hline
\proprow\textbf{Liu et al.} & \textbf{2025} & \textbf{Energy Economics} & \textbf{Carbon (China)} &
  \textbf{VMD--ApEn--LASSO} & \textbf{RMSE, MAE,} \\
\proprow & & & & \textbf{ARIMA/LSTM + Trading} & \textbf{MAPE, DA} \\
\end{tabular}
\renewcommand{\arraystretch}{1}\normalsize
\vspace{8pt}

\subsection{Data Sources and Variables}

The dataset covers 951 weekly observations (January 2000 -- March 2026). The target
is MCX Silver continuous futures price (INR/kg). Ten predictor series are constructed
from nine external sources. Table~2 lists the complete variable inventory used in this
study.

\vspace{4pt}
\noindent\textbf{Table 2.}~The parameter settings in this study.
\vspace{3pt}

\renewcommand{\arraystretch}{1.2}

\noindent\small
\begin{tabular}{L{3.5cm} L{4.2cm} L{5.0cm}}
\rowcolor{HdrBg}\hdrrow{Model / Component} & \hdrrow{Parameter} & \hdrrow{Value / Setting} \\
\hline
ES                       & Smoothing coeff.\ ($\alpha$)         & $\alpha=0.2$ (MLE-optimised) \\
ES                       & Trend component                      & Additive \\
ARIMA                    & Order $(p,d,q)$                      & Selected by AIC grid search \\
ARIMA                    & Differencing ($d$)                   & 1 (confirmed by ADF test) \\
ARIMA                    & Information criterion                & AIC \\
SVR                      & Kernel                               & RBF \\
SVR                      & Regularisation ($C$)                 & 100 \\
SVR                      & Epsilon ($\varepsilon$)              & 0.1 \\
SVR                      & Gamma ($\gamma$)                     & Scale ($1/n_\text{feat} \times \sigma^2_X$) \\
RF                       & Number of trees                      & 100 \\
RF                       & Max features                         & $\sqrt{n_\text{features}}$ \\
RF                       & Min samples split                    & 5 \\
MLP                      & Hidden layers                        & 2 layers: [64, 32] \\
MLP                      & Activation / Dropout                 & ReLU / 0.2 \\
MLP                      & Learning rate / Epochs               & 0.001 (Adam) / 200 (patience=20) \\
ELM                      & Hidden neurons / Activation          & 500 / Sigmoid \\
LSTM (Single)            & Hidden units (layers)                & $64\to64$ \\
LSTM (Single)            & Dropout / lr / Batch                 & 0.2 / 0.001 / 32 (patience=15) \\
LSTM (Single)            & Look-back window                     & 4 weeks \\
LSTM (VMD component)     & Hidden units (layers)                & $64\to64$ \\
LSTM (VMD component)     & Dropout / lr / Batch                 & 0.2 / 0.001 / 16 (patience=15) \\
LASSO                    & Alpha ($\lambda$)                    & CV-selected (5-fold) \\
LASSO                    & Candidate predictors                 & 8 external $+$ 5 AR lags \\
VMD                      & Modes ($K$) / $\alpha$ / tol         & 9 / 2000 / $10^{-7}$ \\
VMD                      & Noise tol.\ ($\tau$) / DC            & 0 / 0 \\
ApEn                     & Embedding dim.\ ($m$) / tol.\ ($r$) & 2 / $0.2\times\sigma$ \\
ApEn                     & Classification rule                  & ApEn(IMF)$<$ApEn(original)$\Rightarrow$ARIMA; else LSTM \\
Trading                  & Interval threshold ($\delta$)        & $\pm3\%$ of predicted price \\
Trading                  & Transaction cost                     & 0.05\% per trade \\
Trading                  & Position sizing                      & Kelly criterion (32.85\%) \\
\end{tabular}
\renewcommand{\arraystretch}{1}\normalsize
\vspace{8pt}

\subsection{Proposed Framework}

Figure~2 shows the structure of the LSTM cell, the core building block of the
high-complexity IMF forecasting stage.

\vspace{4pt}
\noindent\includegraphics[width=0.82\linewidth]{fig2_lstm_architecture.png}
\captionof{figure}{The structure of LSTM. The cell incorporates three gating mechanisms
(input, forget, output) and a memory cell that enables the network to capture long-range
temporal dependencies across the silver price series.}
\label{fig:lstm}
\vspace{8pt}

Figure~3 presents the complete research framework of this study encompassing five
main stages: data collection, data processing, point forecasting, interval
forecasting, and trading simulation.

\vspace{4pt}
\noindent\includegraphics[width=\linewidth]{fig3_research_framework.png}
\captionof{figure}{The framework of this study. The five-stage pipeline proceeds from
multi-source data collection through VMD decomposition, ApEn complexity classification,
LASSO feature selection, adaptive ARIMA/LSTM forecasting, and a four-scheme
interval-constrained trading simulation.}
\label{fig:framework}
\vspace{8pt}

\subsubsection{Variational Mode Decomposition (VMD)}

VMD decomposes the silver price series $s(t)$ into $K=9$ IMFs by solving:
\begin{equation}
\min_{\{u_k\},\{\omega_k\}}\sum_{k=1}^{K}
\Bigl\|\partial_t\!\bigl[(\delta(t)+\tfrac{j}{\pi t})*u_k(t)\bigr]
e^{-j\omega_k t}\Bigr\|_2^2
\quad\text{s.t.}\quad\sum_{k=1}^{K}u_k(t)=s(t)
\label{eq:vmd}
\end{equation}
$K=9$ is estimated as $\lfloor\log_2 N\rfloor$ clamped to $[6,10]$. IMFs are sorted
by ascending dominant frequency. Parameters: $\alpha=2000$, $\tau=0$, DC$=0$,
init$=1$, tol$=10^{-7}$.

\subsubsection{Approximate Entropy Classification}

\begin{equation}
\mathrm{ApEn}(m,r,N)=\Phi^{m}(r)-\Phi^{m+1}(r),\quad
\Phi^m(r)=\tfrac{1}{N-m+1}\sum_{i=1}^{N-m+1}\ln C_i^m(r)
\label{eq:apen}
\end{equation}
Parameters: $m=2$, $r=0.2\times\mathrm{std}(u_k)$. IMFs with ApEn \emph{below} the
original series threshold $\to$ ARIMA; above $\to$ LSTM.

\subsubsection{LASSO Feature Selection}

\begin{equation}
\min_{\bm{\beta}}\Bigl\{\tfrac{1}{2n}\|\mathbf{y}-\mathbf{X}\bm{\beta}\|_2^2
+\alpha\|\bm{\beta}\|_1\Bigr\}
\label{eq:lasso}
\end{equation}
Regularisation $\alpha$ selected by 5-fold \texttt{LassoCV}. Candidate features: five
AR lags of MCX Silver plus nine lagged external variables (14 total per IMF).

\subsubsection{Diebold--Mariano Test}

\begin{equation}
DM=\frac{-\bar{d}}{\sqrt{\widehat{\mathrm{Var}}(d)/n}}
\cdot\sqrt{\frac{n+1-2h+h(h-1)/n}{n}},\quad d_t=e_{1t}^2-e_{2t}^2
\label{eq:dm}
\end{equation}
Positive DM $\Rightarrow$ row model more accurate. HLN correction applied.

\subsubsection{Trading Strategy}

Four schemes with $TC=0.05\%$ per-trade cost. Interval constraint
$I_t=\mathbf{1}[y_t\in[\hat{L}_{t+1},\hat{H}_{t+1}]]$, where
$[\hat{L}_t,\hat{H}_t]=\hat{y}_t\pm3\%$:
\begin{equation}
\widetilde{\mathrm{Signal}}_t=I_t\cdot\mathrm{Signal}_t
\label{eq:constraint}
\end{equation}
\textbf{Scheme~1}: long/short on predicted direction.
\textbf{Scheme~1$'$}: Scheme~1 $\times\,I_t$.
\textbf{Scheme~2}: trade only if $|\Delta\hat{y}_t/\hat{y}_{t-1}|>TC$.
\textbf{Scheme~2$'$}: Scheme~2 $\times\,I_t$.
Sharpe: $(\mu_r/\sigma_r)\times\sqrt{52}$ (annualised).

% ============================================================
\section{Empirical Analysis and Discussion}
% ============================================================

\subsection{Data Collection and Processing}

Figure~4 shows the MCX Silver weekly price series with the train/test boundary.

\vspace{4pt}
\noindent\includegraphics[width=\linewidth]{fig4_silver_price_split.png}
\captionof{figure}{The silver price time series. Green = training set (843 weeks,
Jan 2000 -- Feb 2024); red = testing set (108 weeks, Mar 2024 -- Mar 2026). The
vertical dashed line marks the train/test split. The test slope is $\approx42\times$
steeper than training, constituting a structural regime change.}
\label{fig:pricesplit}
\vspace{8pt}

Table~3 lists all indicators used in this study.

\vspace{4pt}
\noindent\textbf{Table 3.}~The indicators used in this study.
\vspace{3pt}

\renewcommand{\arraystretch}{1.2}

\noindent\small
\begin{tabular}{L{1.7cm} L{3.3cm} L{2.2cm} L{3.0cm} L{1.8cm} L{1.6cm}}
\rowcolor{HdrBg}\hdrrow{Variable} & \hdrrow{Full Name} & \hdrrow{Category} & \hdrrow{Source} & \hdrrow{Frequency} & \hdrrow{Unit} \\
\hline
\texttt{mcx\_silver} & MCX Silver Continuous Futures & Target       & MCX, Bloomberg           & Weekly       & INR/kg \\
\texttt{gold\_usd}   & International Gold Spot Price & Precious Metal& LBMA / Yahoo Finance     & Weekly       & USD/oz \\
\texttt{brent}       & Brent Crude Oil Spot          & Energy        & ICE / Bloomberg          & Weekly       & USD/bbl \\
\texttt{usdinr}      & USD/INR Exchange Rate         & Macro/Currency& RBI / Bloomberg          & Weekly       & INR/USD \\
\texttt{nifty50}     & NIFTY 50 Index (NSE)          & Equity        & NSE India                & Weekly       & Index pts \\
\texttt{vix\_india}  & India VIX                     & Risk/Fear     & NSE India                & Weekly       & Index \\
\texttt{mcx\_gold}   & MCX Gold Continuous Futures   & Domestic PM   & MCX                      & Weekly       & INR/10g \\
\texttt{geo\_risk}   & Geopolitical Risk Index (GPR) & Macro/Geo     & Caldara \& Iacoviello    & Weekly       & Index \\
\texttt{trends\_raw} & Google Trends (12 keywords)   & NLP/Attention & Google Trends API        & Weekly       & 0--100 \\
\texttt{india\_epu}  & India EPU Index               & Policy Risk   & Baker, Bloom \& Davis    & Monthly$^*$  & Index \\
\end{tabular}
\renewcommand{\arraystretch}{1}\normalsize

\noindent\footnotesize $^*$Monthly EPU resampled to weekly via forward-fill.\normalsize
\vspace{8pt}

Figure~5 shows the wordcloud of collected news and search-query terms related to
silver topics in the Indian market, confirming the relevance of the Google Trends
keywords used as a public-attention proxy.

\vspace{4pt}
\noindent\includegraphics[width=0.88\linewidth]{fig5_wordcloud_silver.png}
\captionof{figure}{Wordcloud graph of collected search-query and news terms related
to silver topics in the Indian market. Dominant terms include ``silver'', ``MCX'',
``gold'', ``commodity'', ``rupee'', and ``investment'', validating the Google Trends
keyword selection.}
\label{fig:wordcloud}
\vspace{8pt}

Figure~6 presents the sentiment/attention scores and distributions extracted from
the Google Trends series over the full sample period.

\vspace{4pt}
\noindent\includegraphics[width=\linewidth]{fig6_sentiment_scores.png}
\captionof{figure}{The sentiment scores and distributions extracted from news
headlines and Google Trends. Left: time series of the investor attention index
(2000--2026). Right: distribution of scores showing positive skewness, with spikes
corresponding to major market events and policy announcements.}
\label{fig:sentiment}
\vspace{8pt}

Table~4 presents the statistical information of all data used in this study.

\vspace{4pt}
\noindent\textbf{Table 4.}~The statistical information of used data.
\vspace{3pt}

\renewcommand{\arraystretch}{1.2}

\noindent\scriptsize\setlength{\tabcolsep}{3.5pt}
\begin{tabular}{L{2.2cm} C{0.55cm} C{1.5cm} C{1.5cm} C{1.4cm} C{1.5cm} C{1.6cm}
                C{1.1cm} C{1.1cm} C{1.2cm} C{1.1cm} C{1.1cm}}
\rowcolor{HdrBg}\hdrrow{Variable} & \hdrrow{N} & \hdrrow{Mean} & \hdrrow{Std Dev} & \hdrrow{Min} & \hdrrow{Median} & \hdrrow{Max} & \hdrrow{Skew} & \hdrrow{Kurt} & \hdrrow{JB} & \hdrrow{JB $p$} & \hdrrow{ADF $p$} \\
\hline
MCX Silver (INR/kg) & 951 & 60,289 & 36,276 & 32,872 & 47,195  & 333,292 & 3.73 & 17.93 & 14,942 & 0.00 & 1.000 \\
Gold (USD/oz)       & 951 & 1,622  & 700    & 717    & 1,388   & 5,231   & 2.39 &  7.14 &  2,926 & 0.00 & 1.000 \\
Brent (USD/bbl)     & 951 & 78.08  & 24.05  & 21.44  & 75.15   & 146.08  & 0.23 & $-$0.78 & 33  & 0.00 & 0.037 \\
USD/INR             & 951 & 65.44  & 13.68  & 39.30  & 66.48   & 93.71   & $-$0.09 & $-$1.02 & 43 & 0.00 & 0.909 \\
NIFTY 50            & 951 & 11,137 & 6,488  & 2,584  & 8,834   & 26,329  & 0.88 & $-$0.42 & 128 & 0.00 & 0.986 \\
India VIX           & 951 & 19.73  & 9.00   & 9.15   & 16.88   & 70.39   & 2.20 &  6.36 &  2,372 & 0.00 & 0.018 \\
MCX Gold (INR/10g)  & 951 & 39,784 & 24,764 & 11,560 & 30,283  & 175,231 & 2.18 &  5.87 &  2,121 & 0.00 & 1.000 \\
Geopolitical Risk   & 951 & 103.53 & 55.81  & 0.00   & 92.60   & 478.55  & 1.87 &  7.09 &  2,550 & 0.00 & 0.044 \\
Google Trends       & 951 & 5.00   & 10.48  & 0.00   & 2.33    & 100.00  & 6.46 & 46.86 & 93,631 & 0.00 & 0.223 \\
India EPU Index     & 946 & 105.30 & 49.27  & 23.35  & 92.65   & 327.49  & 1.14 &  1.87 &    344 & 0.00 & 0.114 \\
\end{tabular}
\renewcommand{\arraystretch}{1}\normalsize\setlength{\tabcolsep}{6pt}
\vspace{8pt}

\subsection{VMD Decomposition and Complexity Assessment}

Figure~7 shows the decomposed IMFs of the MCX Silver price series.

\vspace{4pt}
\noindent\includegraphics[width=\linewidth]{fig7_imf_decomposition.png}
\captionof{figure}{The decomposed IMFs of the silver price series. Green = training
portion; red = testing portion. IMF1 (top) captures the dominant long-term secular
trend (variance ratio 71.4\%); IMF9 (bottom) captures high-frequency noise.}
\label{fig:imfs}
\vspace{8pt}

Table~5 presents detailed frequency, period, variance, and correlation statistics for
each decomposed IMF.

\vspace{4pt}
\noindent\textbf{Table 5.}~More information on the decomposed IMFs.
\vspace{3pt}

\renewcommand{\arraystretch}{1.25}

\noindent\begin{tabular}{C{1.0cm} C{2.4cm} C{2.6cm} C{2.4cm} C{2.6cm} C{1.8cm}}
\rowcolor{HdrBg}\hdrrow{Mode} & \hdrrow{Frequency} & \hdrrow{Period (weeks)} & \hdrrow{Variance Ratio} & \hdrrow{Correlation} & \hdrrow{Model} \\
\hline
IMF1 & 0.0011 & 946.00 & 0.7137 & 0.9074 & ARIMA \\
IMF2 & 0.0063 & 157.67 & 0.1471 & 0.5255 & ARIMA \\
IMF3 & 0.0581 &  17.20 & 0.0065 & 0.1332 & LSTM  \\
IMF4 & 0.1015 &   9.85 & 0.0034 & 0.0880 & LSTM  \\
IMF5 & 0.1924 &   5.20 & 0.0013 & 0.0569 & LSTM  \\
IMF6 & 0.2918 &   3.43 & 0.0002 & 0.0264 & LSTM  \\
IMF7 & 0.3372 &   2.97 & 0.0001 & 0.0226 & LSTM  \\
IMF8 & 0.4249 &   2.35 & 0.0004 & 0.0293 & LSTM  \\
IMF9 & 0.4493 &   2.23 & 0.0004 & 0.0301 & LSTM  \\
\end{tabular}
\renewcommand{\arraystretch}{1}
\vspace{8pt}

IMF1 ($\approx$18-year period) alone accounts for 71.4\% of total variance with
correlation 0.907 to the original series. IMF2 ($\approx$3-year supercycle)
contributes 14.7\%. Together they explain 86.1\% of price dynamics. IMFs 3--9 capture
progressively shorter cycles from 17.2 weeks down to 2.23 weeks.

Figure~8 shows the approximate entropy values used to classify each IMF.

\vspace{4pt}
\noindent\includegraphics[width=0.90\linewidth]{fig8_approximate_entropy.png}
\captionof{figure}{The approximate entropy values of decomposed IMFs. Blue bars =
low complexity (ApEn below original series threshold, assigned to ARIMA); red bars =
high complexity (assigned to LSTM). Dashed line = ApEn of the original silver price
series (classification threshold).}
\label{fig:entropy}
\vspace{8pt}

\subsection{Adaptive Feature Selection}

Table~6 presents the selected indicators of decomposed IMFs using LASSO regression,
including the cross-validated regularisation parameter $\alpha$ and number of features
selected $N$.

\vspace{4pt}
\noindent\textbf{Table 6.}~The selected indicators of decomposed IMFs using LASSO
regression. $\checkmark$ = retained (non-zero coefficient). L1--L5 =
\texttt{mcx\_silver\_lag\_1} through \texttt{lag\_5}. External (lag-1): GD = gold\_usd,
BR = brent, FX = usdinr, NF = nifty50, VX = vix\_india, MG = mcx\_gold,
GP = geo\_risk, TR = trends\_raw, EP = india\_epu.
\vspace{3pt}

\renewcommand{\arraystretch}{1.2}

\noindent\scriptsize\setlength{\tabcolsep}{3pt}
\begin{tabular}{C{0.7cm} C{0.9cm} C{1.35cm} C{0.4cm}
  C{0.42cm}C{0.42cm}C{0.42cm}C{0.42cm}C{0.42cm}
  C{0.42cm}C{0.42cm}C{0.42cm}C{0.42cm}C{0.42cm}
  C{0.42cm}C{0.42cm}C{0.42cm}C{0.42cm}}
\rowcolor{HdrBg}\hdrrow{IMF} & \hdrrow{Type} & \hdrrow{$\alpha$} & \hdrrow{N} & \hdrrow{L1}&\hdrrow{L2}&\hdrrow{L3}&\hdrrow{L4}&\hdrrow{L5}& \hdrrow{GD}&\hdrrow{BR}&\hdrrow{FX}&\hdrrow{NF}&\hdrrow{VX}& \hdrrow{MG}&\hdrrow{GP}&\hdrrow{TR}&\hdrrow{EP} \\
\hline
IMF1 & Low  & 22.119   & 11 & \cmark&\cmark&     &     &\cmark&\cmark&\cmark&\cmark&\cmark&\cmark&\cmark&\cmark&\cmark& \\
IMF2 & Low  & 1013.630 &  1 & \cmark&     &     &     &     &     &     &     &     &     &     &     &     & \\
IMF3 & High & 3.699    & 13 & \cmark&\cmark&\cmark&     &\cmark&\cmark&\cmark&\cmark&\cmark&\cmark&\cmark&\cmark&\cmark&\cmark \\
IMF4 & High & 3.620    & 11 & \cmark&     &\cmark&\cmark&\cmark&\cmark&\cmark&\cmark&\cmark&\cmark&     &     &\cmark&\cmark \\
IMF5 & High & 0.861    & 11 & \cmark&\cmark&\cmark&     &\cmark&\cmark&\cmark&     &\cmark&\cmark&     &\cmark&\cmark&\cmark \\
IMF6 & High & 0.185    & 14 & \cmark&\cmark&\cmark&\cmark&\cmark&\cmark&\cmark&\cmark&\cmark&\cmark&\cmark&\cmark&\cmark&\cmark \\
IMF7 & High & 0.077    & 14 & \cmark&\cmark&\cmark&\cmark&\cmark&\cmark&\cmark&\cmark&\cmark&\cmark&\cmark&\cmark&\cmark&\cmark \\
IMF8 & High & 0.003    & 14 & \cmark&\cmark&\cmark&\cmark&\cmark&\cmark&\cmark&\cmark&\cmark&\cmark&\cmark&\cmark&\cmark&\cmark \\
IMF9 & High & 0.324    & 11 & \cmark&\cmark&\cmark&\cmark&\cmark&     &\cmark&\cmark&\cmark&\cmark&     &     &\cmark&\cmark \\
\end{tabular}
\renewcommand{\arraystretch}{1}\normalsize\setlength{\tabcolsep}{6pt}
\vspace{6pt}

Key patterns: (1)~\texttt{mcx\_silver\_lag\_1} selected in all 9~IMFs, confirming
strong short-run autocorrelation; (2)~India EPU selected in 8/9~IMFs---the most
pervasive external driver; (3)~IMF2 selects only lag-1 (zero external features),
consistent with a self-reinforcing 3-year supercycle; (4)~IMF6--IMF8 select all
14~features; (5)~high-complexity IMFs depend more on historical lags than
low-complexity IMFs, consistent with the findings of Liu et al.\ (2025).

\subsection{Point Forecasting Results}

Table~7 reports forecasting errors for all single-model benchmarks on the 108-week
test set.

\vspace{4pt}
\noindent\textbf{Table 7.}~Forecasting errors on the testing sets for single models.
DA = directional accuracy; DA $p$-val tests H$_0$: DA$=50\%$ (one-sided binomial).
$^{***}p<0.01$, $^{**}p<0.05$, $^*p<0.10$.
\vspace{3pt}

\renewcommand{\arraystretch}{1.25}

\noindent\begin{tabular}{L{3.4cm} R{2.2cm} R{2.2cm} R{1.6cm} R{1.7cm} R{1.3cm} R{1.6cm}}
\rowcolor{HdrBg}\hdrrow{Model} & \hdrrow{RMSE} & \hdrrow{MAE} & \hdrrow{MAPE (\%)} & \hdrrow{sMAPE (\%)} & \hdrrow{DA (\%)} & \hdrrow{DA $p$-val} \\
\hline
ES                  & 44,434.43  & 24,157.08  & 13.77 & 15.31 & 54.90 & 0.1865 \\
ARIMA               & 71,519.30  & 43,552.17  & 24.95 & 32.41 & 31.37 & 1.0000 \\
SVR                 & 100,694.08 & 83,094.32  & 59.82 & 87.45 & 31.37 & 1.0000 \\
RF                  & 69,937.99  & 40,934.56  & 22.65 & 29.53 & 39.22 & 0.9889 \\
MLP                 & 123,461.59 & 113,941.25 & 89.33 & 161.57& 31.37 & 1.0000 \\
ELM                 & 97,421.38  & 76,407.25  & 54.53 & 80.10 & 31.37 & 1.0000 \\
LSTM (single)       & 140,545.95 & 128,552.85 & 99.99 & 199.96& 31.37 & 1.0000 \\
Naive (Random Walk) & 11,219.77  &  5,247.51  &  3.26 &  3.26 &  0.00 & 1.0000 \\
\end{tabular}
\renewcommand{\arraystretch}{1}
\vspace{8pt}

Table~8 reports forecasting errors for all decomposition-based models.

\vspace{4pt}
\noindent\textbf{Table 8.}~Forecasting errors on the testing sets for the models
with decomposition techniques. Proposed row highlighted green.
\vspace{3pt}

\renewcommand{\arraystretch}{1.25}

\noindent\begin{tabular}{L{3.4cm} R{2.2cm} R{2.2cm} R{1.6cm} R{1.7cm} R{1.3cm} R{1.6cm}}
\rowcolor{HdrBg}\hdrrow{Model} & \hdrrow{RMSE} & \hdrrow{MAE} & \hdrrow{MAPE (\%)} & \hdrrow{sMAPE (\%)} & \hdrrow{DA (\%)} & \hdrrow{DA $p$-val} \\
\hline
VMD-ARIMA     & 64,509.63  & 35,049.78  & 18.32 & 23.60 & 50.00 & 0.5394 \\
VMD-LSTM      & 140,548.71 & 128,554.03 & 99.99 & 199.96& 31.37 & 1.0000 \\
CEEMDAN-ARIMA & 64,372.43  & 35,055.17  & 18.36 & 23.62 & 50.98 & 0.4606 \\
CEEMDAN-LSTM  & 140,508.69 & 128,503.80 & 99.93 & 199.01& 31.37 & 1.0000 \\
\proprow\textbf{Proposed} &
  \textbf{63,708.61} & \textbf{34,343.79} & \textbf{17.88} &
  \textbf{22.94} & \textbf{46.08} & \textbf{0.8135} \\
\end{tabular}
\renewcommand{\arraystretch}{1}
\vspace{8pt}

Figure~9 shows the forecasting results using single forecasting models on the test set.

\vspace{4pt}
\noindent\includegraphics[width=\linewidth]{fig9_single_model_forecasts.png}
\captionof{figure}{The forecasting results using single forecasting models. Dark line
= actual MCX Silver price; coloured dashed lines = model predictions. High RMSE
values reflect the unprecedented 2024--2026 bull market structural break.}
\label{fig:singleforecasts}
\vspace{8pt}

Figure~10 presents bar plots of forecasting error metrics across all models.

\vspace{4pt}
\noindent\includegraphics[width=\linewidth]{fig10_error_barplots.png}
\captionof{figure}{Bar plots of forecasting error metrics across models. Blue =
single models; red = decomposition benchmarks; green = Proposed. The Proposed model
achieves the lowest RMSE (63,709) and MAE (34,344) among all decomposition models.}
\label{fig:barplots}
\vspace{8pt}

The Proposed model achieves RMSE~63,709 and MAE~34,344 (best among decomposition
models), outperforming VMD-ARIMA by 0.9\% and CEEMDAN-ARIMA by 1.0\%. Critically,
VMD-LSTM and CEEMDAN-LSTM catastrophically underperform (RMSE~$\approx140{,}000$),
confirming that applying LSTM uniformly across all IMFs---including the smooth secular
trend---causes severe overfitting.

\subsection{Statistical Validation: Diebold--Mariano Test}

Table~9 presents the full $12\times12$ DM test results.

\vspace{4pt}
\noindent\textbf{Table 9.}~The results of DM-test. Positive statistic = row model more
accurate. $^{***}p<0.01$, $^{**}p<0.05$, $^*p<0.10$. Proposed row highlighted green.
\vspace{3pt}

\renewcommand{\arraystretch}{1.15}

\noindent\scriptsize\setlength{\tabcolsep}{3pt}
\begin{tabular}{L{2.1cm} C{1.1cm} C{1.1cm} C{1.1cm} C{1.1cm} C{1.1cm} C{1.1cm}
                C{1.1cm} C{1.3cm} C{1.3cm} C{1.5cm} C{1.5cm}}
\rowcolor{HdrBg}\hdrrow{Row vs Col.} & \hdrrow{ES} & \hdrrow{ARIMA} & \hdrrow{SVR} & \hdrrow{RF} & \hdrrow{MLP} & \hdrrow{ELM} & \hdrrow{LSTM} & \hdrrow{VMD-AR} & \hdrrow{VMD-LS} & \hdrrow{CEEM-AR} & \hdrrow{CEEM-LS} \\
\hline
ES            & ---             & $5.31^{***}$    & $7.87^{***}$    & $5.19^{***}$    & $13.22^{***}$   & $6.20^{***}$    & $11.48^{***}$   & $4.86^{***}$    & $11.48^{***}$   & $4.88^{***}$    & $11.49^{***}$ \\
ARIMA         & $-5.31^{***}$   & ---             & $11.17^{***}$   & $-7.46^{***}$   & $23.16^{***}$   & $6.43^{***}$    & $15.21^{***}$   & $-6.67^{***}$   & $15.21^{***}$   & $-6.52^{***}$   & $15.23^{***}$ \\
SVR           & $-7.87^{***}$   & $-11.17^{***}$  & ---             & $-10.94^{***}$  & $36.75^{***}$   & $-1.47$         & $18.77^{***}$   & $-10.08^{***}$  & $18.76^{***}$   & $-10.02^{***}$  & $18.79^{***}$ \\
RF            & $-5.19^{***}$   & $7.46^{***}$    & $10.94^{***}$   & ---             & $22.22^{***}$   & $6.54^{***}$    & $14.98^{***}$   & $-6.44^{***}$   & $14.98^{***}$   & $-6.24^{***}$   & $14.99^{***}$ \\
MLP           & $-13.22^{***}$  & $-23.16^{***}$  & $-36.75^{***}$  & $-22.22^{***}$  & ---             & $-10.99^{***}$  & $8.02^{***}$    & $-19.29^{***}$  & $8.02^{***}$    & $-19.15^{***}$  & $8.03^{***}$  \\
ELM           & $-6.20^{***}$   & $-6.43^{***}$   & $1.47$          & $-6.54^{***}$   & $10.99^{***}$   & ---             & $16.65^{***}$   & $-6.69^{***}$   & $16.65^{***}$   & $-6.65^{***}$   & $16.61^{***}$ \\
LSTM          & $-11.48^{***}$  & $-15.21^{***}$  & $-18.77^{***}$  & $-14.98^{***}$  & $-8.02^{***}$   & $-16.65^{***}$  & ---             & $-14.11^{***}$  & $2.66^{***}$    & $-14.06^{***}$  & $-0.65$       \\
VMD-ARIMA     & $-4.86^{***}$   & $6.67^{***}$    & $10.08^{***}$   & $6.44^{***}$    & $19.29^{***}$   & $6.69^{***}$    & $14.11^{***}$   & ---             & $14.11^{***}$   & $-1.18$         & $14.12^{***}$ \\
VMD-LSTM      & $-11.48^{***}$  & $-15.21^{***}$  & $-18.76^{***}$  & $-14.98^{***}$  & $-8.02^{***}$   & $-16.65^{***}$  & $-2.66^{***}$   & $-14.11^{***}$  & ---             & $-14.06^{***}$  & $-0.70$       \\
CEEMDAN-ARIMA & $-4.88^{***}$   & $6.52^{***}$    & $10.02^{***}$   & $6.24^{***}$    & $19.15^{***}$   & $6.65^{***}$    & $14.06^{***}$   & $1.18$          & $14.06^{***}$   & ---             & $14.07^{***}$ \\
CEEMDAN-LSTM  & $-11.49^{***}$  & $-15.23^{***}$  & $-18.79^{***}$  & $-14.99^{***}$  & $-8.03^{***}$   & $-16.61^{***}$  & $0.65$          & $-14.12^{***}$  & $0.70$          & $-14.07^{***}$  & ---           \\
\proprow\textbf{Proposed} &
  $\mathbf{-4.81^{***}}$ & $\mathbf{6.60^{***}}$ & $\mathbf{9.97^{***}}$ &
  $\mathbf{6.39^{***}}$ & $\mathbf{18.94^{***}}$ & $\mathbf{6.70^{***}}$ &
  $\mathbf{13.99^{***}}$ & $\mathbf{5.89^{***}}$ & $\mathbf{13.99^{***}}$ &
  $\mathbf{4.42^{***}}$ & $\mathbf{14.00^{***}}$ \\
\end{tabular}
\renewcommand{\arraystretch}{1}\normalsize\setlength{\tabcolsep}{6pt}
\vspace{6pt}

\begin{insightbox}[title=\textbf{KEY FINDING}]
VMD-LSTM scores $DM=-14.11^{***}$ against VMD-ARIMA, meaning applying LSTM uniformly
across all IMFs is \emph{significantly worse} than uniform ARIMA. The Proposed model
beats VMD-ARIMA at $DM=5.89^{***}$, confirming that ApEn-guided complexity assignment
is the critical methodological contribution.
\end{insightbox}

\subsection{Interval Forecasting}

Figure~11 presents the interval forecasting results.

\vspace{4pt}
\noindent\includegraphics[width=\linewidth]{fig11_interval_forecasts.png}
\captionof{figure}{The interval forecasting results. Gray shaded band = predicted
$\pm3\%$ interval; green line = point prediction (Proposed model); red dashed =
actual MCX Silver price. The coverage ratio (CR$=0.2816$) reflects the structural
break: the band, calibrated on training volatility, is insufficient for the
unprecedented 2024--2026 bull market.}
\label{fig:interval}
\vspace{8pt}

Table~10 presents interval forecasting errors on the testing set using the iMLP model.

\vspace{4pt}
\noindent\textbf{Table 10.}~Interval forecasting errors on the testing sets using
the iMLP model. U = Theil statistic; ARV = Average Relative Variance; RMSDE = Root
Mean Square Distance Error; CR = Coverage Ratio.
\vspace{3pt}

\renewcommand{\arraystretch}{1.25}

\noindent\begin{tabular}{L{4.5cm} C{2.0cm} C{2.0cm} C{2.6cm} C{2.0cm}}
\rowcolor{HdrBg}\hdrrow{Model} & \hdrrow{U} & \hdrrow{ARV} & \hdrrow{RMSDE} & \hdrrow{CR} \\
\hline
Silver (Proposed) & 0.1014 & 1.2577 & 63,708.61 & 0.2816 \\
\end{tabular}
\renewcommand{\arraystretch}{1}
\vspace{8pt}

The coverage ratio of 28.2\% is below the nominal 95\% target---a structural
consequence of the test-period regime change, where prices surged far beyond any band
calibrated on training-period volatility (ARV$>1$).

\subsection{Trading Strategy Results}

Figure~12 provides a graph illustration of the proposed trading strategy mechanism.

\vspace{4pt}
\noindent\includegraphics[width=\linewidth]{fig12_trading_strategy_illustration.png}
\captionof{figure}{Graph illustration of the proposed trading strategy. Green band =
predicted interval; purple = point prediction; red dashed = actual price; red circles
= weeks where point forecast falls outside the predicted interval (constraint active,
trade suppressed).}
\label{fig:strategy}
\vspace{8pt}

Table~11 reports the performance of both trading schemes and their interval-constrained
versions based on decomposition model forecasts.

\vspace{4pt}
\noindent\textbf{Table 11.}~Performance of two trading schemes and their optimised
versions with interval constraints based on decomposition model forecasts. Enhancement
in cumulative return shown in parentheses (bold). Most profitable per scheme
underlined. Asterisk (*) = loss in unconstrained scheme.
\vspace{3pt}

\renewcommand{\arraystretch}{1.18}

\noindent\small
\begin{tabular}{L{1.5cm} L{2.5cm} R{2.5cm} R{1.6cm} R{1.6cm} R{1.4cm} R{1.3cm}}
\rowcolor{HdrBg}\hdrrow{Scheme} & \hdrrow{Model} & \hdrrow{Cum.\ Return (\%)} & \hdrrow{Avg Ret (\%)} & \hdrrow{Max DD (\%)} & \hdrrow{Sharpe} & \hdrrow{N Trades} \\
\hline
\multirow{5}{=}{Scheme 1}
  & VMD-ARIMA      &   7.5182    & 0.1994   & 50.6279 & 0.2945   & 102 \\
  & VMD-LSTM       & $-$52.4623  & $-$0.6036& 64.6692 & $-$0.8962& 102 \\
  & CEEMDAN-ARIMA  &   8.4349    & 0.2056   & 45.0913 & 0.3037   & 102 \\
  & CEEMDAN-LSTM   & $-$70.6649  & $-$1.0820& 76.5415 & $-$1.6331& 102 \\
  & \underline{\textbf{Proposed}} & \underline{\textbf{187.1550}} & \textbf{1.1556} & \textbf{27.7176} & \textbf{1.7591} & 102 \\
\hline
\multirow{5}{=}{Scheme 1$'$}
  & VMD-ARIMA      & 52.1144 (\textbf{+44.60})   & 0.4229   &  4.3968 & 2.0493   & 30 \\
  & VMD-LSTM       & 30.0091 (\textbf{+82.47})   & 0.2691   &  4.5484 & 1.2693   & 30 \\
  & CEEMDAN-ARIMA  & 23.6603 (\textbf{+15.23})   & 0.2202   &  8.1908 & 1.0326   & 30 \\
  & CEEMDAN-LSTM   & $-$21.8179 (\textbf{+48.85})& $-$0.2289& 29.9639 & $-$1.0639& 30 \\
  & \textbf{Proposed} & \textbf{15.1784} (\textbf{$-$171.98}) & 0.1505 & \textbf{10.8707} & \textbf{0.7010} & 30 \\
\hline
\multirow{5}{=}{Scheme 2}
  & VMD-ARIMA      &  14.7211    & 0.2463   & 49.8568 & 0.3917   &  92 \\
  & VMD-LSTM       & $-$52.4623  & $-$0.6036& 64.6692 & $-$0.8962& 102 \\
  & CEEMDAN-ARIMA  &   8.4349    & 0.2056   & 45.0913 & 0.3037   & 102 \\
  & CEEMDAN-LSTM   & $-$70.6649  & $-$1.0820& 76.5415 & $-$1.6331& 102 \\
  & \underline{\textbf{Proposed}} & \underline{\textbf{23.1753}} & \textbf{0.2234} & \textbf{10.9347} & \textbf{0.8311} & 56 \\
\hline
\multirow{5}{=}{Scheme 2$'$}
  & VMD-ARIMA      & 43.7354 (\textbf{+29.01})   & 0.3651   &  3.3695 & 1.9698   & 27 \\
  & VMD-LSTM       & 30.0091 (\textbf{+82.47})   & 0.2691   &  4.5484 & 1.2693   & 30 \\
  & CEEMDAN-ARIMA  & 23.6603 (\textbf{+15.23})   & 0.2202   &  8.1908 & 1.0326   & 30 \\
  & CEEMDAN-LSTM   & $-$21.8179 (\textbf{+48.85})& $-$0.2289& 29.9639 & $-$1.0639& 30 \\
  & \textbf{Proposed} & \textbf{7.7751} (\textbf{$-$15.40}) & 0.0832 & \textbf{10.8707} & \textbf{0.4295} & 29 \\
\end{tabular}
\renewcommand{\arraystretch}{1}\normalsize
\vspace{8pt}

Table~12 reports the performance of the schemes based on single-model forecasts.

\vspace{4pt}
\noindent\textbf{Table 12.}~Performance of two trading schemes and their optimised
versions with interval constraints based on single model forecasts. Enhancement in
cumulative return shown in parentheses. Asterisk (*) = loss in unconstrained scheme.
\vspace{3pt}

\renewcommand{\arraystretch}{1.18}

\noindent\small
\begin{tabular}{L{1.5cm} L{1.8cm} R{2.6cm} R{1.6cm} R{1.6cm} R{1.4cm} R{1.3cm}}
\rowcolor{HdrBg}\hdrrow{Scheme} & \hdrrow{Model} & \hdrrow{Cum.\ Return (\%)} & \hdrrow{Avg Ret (\%)} & \hdrrow{Max DD (\%)} & \hdrrow{Sharpe} & \hdrrow{N Trades} \\
\hline
\multirow{7}{=}{Scheme 1}
  & ES    & \underline{187.1550} & 1.1556   & 27.7176 & 1.7591   & 102 \\
  & ARIMA & $-$7.9093$^*$        & $-$0.0722& 15.3129 & $-$0.4015&  15 \\
  & SVR   & $-$15.7139$^*$       & $-$0.1191& 38.9948 & $-$0.2744&  86 \\
  & RF    & 81.8759              & 0.6803   & 14.7327 & 1.1125   &  94 \\
  & MLP   & 17.4717              & 0.2814   & 27.7560 & 0.4161   & 102 \\
  & ELM   & 25.6681              & 0.3412   & 29.7820 & 0.5049   & 102 \\
  & LSTM  & $-$38.4903$^*$       & $-$0.3572& 50.8125 & $-$0.5280& 102 \\
\hline
\multirow{7}{=}{Scheme 1$'$}
  & ES    & 15.1784 (\textbf{$-$171.98}) & 0.1505   & 10.8707 & 0.7010   & 30 \\
  & ARIMA & $-$5.6036 (\textbf{+2.31})   & $-$0.0529&  9.2022 & $-$0.4529&  8 \\
  & SVR   & $-$23.2079 (\textbf{$-$7.49})& $-$0.2464& 24.1216 & $-$1.1470& 30 \\
  & RF    & 33.2624 (\textbf{$-$48.61})  & 0.2930   &  4.4820 & 1.4021   & 29 \\
  & MLP   & $-$21.7091 (\textbf{$-$39.18})& $-$0.2273& 23.5457 & $-$1.0566& 30 \\
  & ELM   & $-$2.8003 (\textbf{$-$28.47})& $-$0.0156& 17.2831 & $-$0.0721& 30 \\
  & LSTM  & $-$23.2731 (\textbf{+15.22}$^*$)& $-$0.2473& 28.2278 & $-$1.1513& 30 \\
\hline
\multirow{7}{=}{Scheme 2}
  & ES    & \underline{187.1550} & 1.1556   & 27.7176 & 1.7591   & 102 \\
  & ARIMA & 0.0000               & 0.0000   &  0.0000 & 0.0000   &   0 \\
  & SVR   & $-$15.0569$^*$       & $-$0.1527& 15.0569 & $-$0.9178&  22 \\
  & RF    & 90.3221              & 0.7201   & 16.1343 & 1.2085   &  81 \\
  & MLP   & 17.4717              & 0.2814   & 27.7560 & 0.4161   & 102 \\
  & ELM   & 15.7207              & 0.2569   & 29.7820 & 0.3857   & 101 \\
  & LSTM  & $-$31.7569$^*$       & $-$0.2588& 45.4279 & $-$0.3877&  89 \\
\hline
\multirow{7}{=}{Scheme 2$'$}
  & ES    & 15.1784 (\textbf{$-$171.98}) & 0.1505   & 10.8707 & 0.7010   & 30 \\
  & ARIMA & 0.0000                       & 0.0000   &  0.0000 & 0.0000   &   0 \\
  & SVR   & $-$13.4854 (\textbf{+1.57}$^*$)& $-$0.1392& 13.4854 & $-$1.3761& 11 \\
  & RF    & 24.1569 (\textbf{$-$66.17})  & 0.2226   &  5.4019 & 1.1073   & 26 \\
  & MLP   & $-$21.7091 (\textbf{$-$39.18})& $-$0.2273& 23.5457 & $-$1.0566& 30 \\
  & ELM   & $-$2.8003 (\textbf{$-$18.52})& $-$0.0156& 17.2831 & $-$0.0721& 30 \\
  & LSTM  & $-$19.5534 (\textbf{+12.20}$^*$)& $-$0.2018& 24.7242 & $-$0.9703& 26 \\
\end{tabular}
\renewcommand{\arraystretch}{1}\normalsize
\vspace{8pt}

Figure~13 presents bar plots of trading evaluation metrics across all models and
schemes.

\vspace{4pt}
\noindent\includegraphics[width=\linewidth]{fig13_trading_evaluation.png}
\captionof{figure}{Bar plots of trading evaluation metrics across models. Left:
cumulative return (\%). Right: annualised Sharpe ratio. The Proposed model achieves
the highest Sharpe (1.759) and cumulative return (187.2\%) under Scheme~1 among
decomposition-based models.}
\label{fig:trading}
\vspace{8pt}

The Proposed model under Scheme~1 achieves 187.2\% cumulative return (Sharpe 1.759,
max drawdown 27.7\%) across 102 trades---the highest among all decomposition models.
Under the interval-constrained Scheme~1$'$, the trade count falls from 102 to 30
(a 70.6\% reduction) and maximum drawdown drops from 27.7\% to 10.9\% ($-60.7\%$).
The Conditional DA for Scheme~1$'$ is 56.67\% (17 correct out of 30 trades;
profit factor 1.62; Kelly criterion 21.69\%), demonstrating genuine directional skill
above the 50\% random benchmark. Among single-model strategies, ARIMA, SVR, and LSTM
generate losses under Scheme~1 (marked $^*$). The interval constraint partially
reverses losses (LSTM: $-38.5\%\to-23.3\%$; SVR Scheme~2: $-15.1\%\to-13.5\%$),
consistent with Liu et al.\ (2025)'s finding that interval constraints can reverse
losses to gains.

\begin{insightbox}[title=\textbf{KEY FINDING}]
The interval constraint (Schemes 1$'$/2$'$) reduces the trade count from 102 to 30
weeks and cuts maximum drawdown by 60.7\% (from 27.72\% to 10.87\%), demonstrating
that interval forecasts serve as a principled risk-management tool even during
structural regime changes. The Conditional DA of 56.67\% on filtered trades
(profit factor 1.62) confirms genuine signal quality beyond random chance.
\end{insightbox}

% ============================================================
\section{Conclusion}
% ============================================================

This paper applies the multimodal VMD--ApEn--LASSO framework of Liu et al.\ (2025)
to MCX Silver futures (INR/kg) over 951 weekly observations (2000--2026), with a
108-week test period covering the unprecedented 2024--2026 Indian silver bull market.

On \textbf{decomposition and complexity}, VMD extracts nine IMFs spanning periods from
2.23 to 946~weeks. IMFs~1--2 account for 86.1\% of variance. ApEn classifies IMFs~1--2
as low complexity (ARIMA) and IMFs~3--9 as high complexity (LSTM). VMD-LSTM
($DM=-14.11^{***}$ vs VMD-ARIMA) confirms that complexity-aware assignment is essential.

On \textbf{feature selection}, India EPU is selected in 8/9 IMFs---the most pervasive
driver. Gold co-movement and USD/INR dominate long-run IMFs; India VIX, Google Trends,
and EPU dominate short-term oscillations. IMF2 is purely autoregressive.

On \textbf{forecasting}, the Proposed model achieves RMSE~63,709~INR/kg (best among
decomposition models). DM tests confirm $p<0.01$ superiority over all 11 benchmarks
except ES ($-4.81^{***}$). Significance vs VMD-ARIMA: $5.89^{***}$.

On \textbf{trading}, Scheme~1 generates 187.2\% return (Sharpe 1.759, max~DD 27.7\%).
Scheme~1$'$ cuts maximum drawdown to 10.9\% with 30 filtered trades (Conditional
DA~56.67\%, profit factor~1.62), providing a risk-managed alternative for conservative
institutional participants.

\medskip\noindent\textbf{Limitations and future research.} The fixed $\pm3\%$ interval
band should be replaced by adaptive calibration (e.g., GARCH-LSTM). Daily MCX data
would improve MAPE during rapid price accelerations. Extensions to MCX Gold, Crude
Oil, and Natural Gas would validate entropy-guided assignment across Indian commodity
markets. Temporal Fusion Transformers and richer features (SEBI announcements, MCX
open interest) could further enhance performance.

\newpage

% ============================================================
\appendix
\section*{Appendix.}
% ============================================================

Supplementary data including all per-IMF ARIMA orders, LASSO coefficient paths,
weekly trading signals, and extended DM test results are available from the authors
upon request. Full pipeline code: \texttt{Crystal-Blaze13/Silver\_RC}.

\newpage

% ============================================================
\section*{References}
% ============================================================

\begin{enumerate}[leftmargin=0pt,label={},itemsep=5pt]

\item Liu, S., Li, M., Yang, K., Wei, Y., \& Wang, S. (2025).
From forecasting to trading: A multimodal-data-driven approach to reversing carbon
market losses. \textit{Energy Economics}, \textit{144}, 108350.
\url{https://doi.org/10.1016/j.eneco.2025.108350}

\end{enumerate}

\end{document}